\documentclass[11pt]{article}
\usepackage[margin=1in]{geometry}
\usepackage[T1]{fontenc}
\usepackage[utf8]{inputenc}
\usepackage{lmodern}
\usepackage{microtype}
\usepackage[table]{xcolor}
\usepackage{amsmath,amssymb,amsfonts,mathtools,bm}
\IfFileExists{physics.sty}{%
  \usepackage{physics}%
}{%
  % Portable subset used by this project when the optional physics package
  % is not installed in the local TeX distribution.
  \NewDocumentCommand{\pdv}{o m m}{%
    \IfNoValueTF{##1}{\frac{\partial ##2}{\partial ##3}}%
    {\frac{\partial^{##1} ##2}{\partial ##3^{##1}}}%
  }
  \NewDocumentCommand{\dv}{o m m}{%
    \IfNoValueTF{##1}{\frac{\mathrm{d} ##2}{\mathrm{d} ##3}}%
    {\frac{\mathrm{d}^{##1} ##2}{\mathrm{d} ##3^{##1}}}%
  }
  \providecommand{\dd}{\mathop{}\!\mathrm{d}}
  \providecommand{\grad}{\nabla}
  \providecommand{\abs}[1]{\left\lvert ##1\right\rvert}
  \providecommand{\norm}[1]{\left\lVert ##1\right\rVert}
  \providecommand{\ket}[1]{\left\lvert ##1\right\rangle}
  \providecommand{\bra}[1]{\left\langle ##1\right\rvert}
}
\usepackage{booktabs,array,longtable,tabularx}
\usepackage{hyperref}
\usepackage{enumitem}
\usepackage{fancyhdr}
\usepackage{titlesec}
\usepackage{tcolorbox}
\usepackage{ragged2e}
\usepackage{setspace}
\IfFileExists{siunitx.sty}{%
  \usepackage{siunitx}%
}{%
  % Minimal text-safe fallback for the small number of \SI and \si calls in
  % this repository. Full siunitx formatting is used when available.
  \providecommand{\si}[1]{\ensuremath{\mathrm{##1}}}
  \providecommand{\SI}[2]{\ensuremath{##1\,\mathrm{##2}}}
}
\hypersetup{colorlinks=true, linkcolor=blue!50!black, urlcolor=blue!50!black, citecolor=blue!50!black}
\definecolor{TitleBlue}{HTML}{163A5F}
\definecolor{AccentTeal}{HTML}{2D6F73}
\definecolor{SoftBlue}{HTML}{EAF3F8}
\definecolor{SoftTeal}{HTML}{EDF8F6}
\definecolor{SoftGray}{HTML}{F5F7FA}
\definecolor{RuleGray}{HTML}{D7DEE8}
\pagestyle{fancy}
\fancyhf{}
\lhead{RadiationEffects\_proj2}
\rhead{Technical Notes}
\cfoot{\thepage}
\setlength{\headheight}{14pt}
\setlength{\parskip}{0.5em}
\setlength{\parindent}{0pt}
\onehalfspacing
\numberwithin{equation}{section}
\titleformat{\section}{\Large\bfseries\color{TitleBlue}}{\thesection}{0.6em}{}
\titleformat{\subsection}{\large\bfseries\color{AccentTeal}}{\thesubsection}{0.6em}{}
\titleformat{\subsubsection}{\normalsize\bfseries\color{TitleBlue}}{\thesubsubsection}{0.6em}{}
\newtcolorbox{overviewbox}{colback=SoftBlue,colframe=TitleBlue,boxrule=0.8pt,arc=2mm,left=3mm,right=3mm,top=2mm,bottom=2mm}
\newtcolorbox{physicsbox}{colback=SoftTeal,colframe=AccentTeal,boxrule=0.8pt,arc=2mm,left=3mm,right=3mm,top=2mm,bottom=2mm}
\newtcolorbox{notebox}{colback=SoftGray,colframe=AccentTeal,boxrule=0.6pt,arc=2mm,left=3mm,right=3mm,top=2mm,bottom=2mm}
\newcommand{\DocTitle}[3]{%
\begin{center}
{\LARGE \bfseries \color{TitleBlue} #1}\\[0.5em]
{\large \color{AccentTeal} #2}\\[0.8em]
{\normalsize #3}
\end{center}
\vspace{0.5em}}
\newenvironment{symtable}{\rowcolors{2}{SoftGray}{white}\begin{longtable}{>{\RaggedRight\arraybackslash}p{0.18\textwidth} >{\RaggedRight\arraybackslash}p{0.25\textwidth} >{\RaggedRight\arraybackslash}p{0.47\textwidth}}\toprule \textbf{Symbol / acronym} & \textbf{Name} & \textbf{Meaning in this document} \\ \midrule\endfirsthead\toprule \textbf{Symbol / acronym} & \textbf{Name} & \textbf{Meaning in this document} \\ \midrule\endhead}{\bottomrule\end{longtable}}


\newcommand{\ProjectTOC}{%
\vspace{0.6em}
{\hypersetup{linkcolor=TitleBlue,urlcolor=TitleBlue,citecolor=TitleBlue}\tableofcontents}
\vspace{1.0em}}

\newcommand{\SymbolsAcronymsSection}{%
\section*{Symbols and acronyms used in this document}%
\addcontentsline{toc}{section}{Symbols and acronyms used in this document}}

\newcommand{\rvec}{\bm{r}}
\newcommand{\qvec}{\bm{q}}
\newcommand{\nvec}{\bm{n}}
\newcommand{\xvec}{\bm{x}}
\newcommand{\kB}{k_{\mathrm{B}}}
\newcommand{\qp}{\mathrm{qp}}
\newcommand{\JJ}{\mathrm{JJ}}
\newcommand{\mfp}{\Lambda}
\allowdisplaybreaks

\renewcommand{\xvec}{\mathbf{x}}
\newcommand{\Jvec}{\mathbf{J}}
\newcommand{\Ddef}{D}
\newcommand{\kann}{k_{\mathrm{ann}}}
\newcommand{\Rrxn}{R}
\newcommand{\Cdef}{C}
\newcommand{\Tfield}{T}
\rhead{Module 3 Physics Note}

\begin{document}

\DocTitle{Module 3: Two-Dimensional Defect Diffusion-Reaction Evolution and Annealing}{Derivation-Focused Physics Note for \texttt{RadiationEffects\_proj2}}{Time evolution of radiation-induced defects in silicon after the initial damage event}

\hrule
\vspace{0.8em}

\begin{overviewbox}
\textbf{Purpose of this note.} This note develops the defect-evolution module that turns an initial radiation damage map into a time-dependent defect landscape inside silicon. The goal is to derive, from conservation of defect number and constitutive transport ideas, the two-dimensional diffusion-reaction equations used to model defect migration, recombination, clustering, dissociation, and annealing. This module provides the evolving damage state that later feeds the electrostatic, thermal, and carrier-transport parts of the project.
\end{overviewbox}

\ProjectTOC


\SymbolsAcronymsSection

\renewcommand{\arraystretch}{1.25}
\begin{longtable}{p{0.22\textwidth} p{0.56\textwidth} p{0.16\textwidth}}
\toprule
\textbf{Symbol} & \textbf{Meaning} & \textbf{Typical units} \\
\midrule
\endhead
$C$ & Defect concentration field & m$^{-3}$ \\
$C_i$ & Concentration of defect species $i$ & m$^{-3}$ \\
$\mathbf{J}_C$ & Defect flux & m$^{-2}$ s$^{-1}$ \\
$D$ & Defect diffusivity & m$^2$/s \\
$D_0$ & Diffusion prefactor, typically $a^2 \nu_D /(2d)$ in a jump-diffusion picture & m$^2$/s \\
$R$ & Net local reaction term & m$^{-3}$ s$^{-1}$ \\
$k_{\mathrm{ann}}$ & First-order annealing rate constant & s$^{-1}$ \\
$k_0$ & Annealing prefactor or attempt rate for the modeled annealing channel & s$^{-1}$ \\
$T$ & Temperature field & K \\
$E_D$ & Effective migration activation energy for defect motion between metastable lattice configurations & J or eV \\
$E_{\mathrm{ann}}$ & Effective activation energy for defect removal or transformation during annealing & J or eV \\
$k_B$ & Boltzmann constant & J/K \\
$x,y$ & Cartesian spatial coordinates & m \\
$t$ & Time & s \\
$L$ & Characteristic length scale & m \\
$\tau_D$ & Diffusion timescale & s \\
$\tau_{\mathrm{ann}}$ & Annealing timescale & s \\
$M(t)$ & Total defect inventory & dimension depends on dimensional reduction \\
$\hat{\mathbf{n}}$ & Outward unit normal vector & dimensionless \\
$\Omega$ & Two-dimensional computational domain & m$^2$ \\
$\Omega_e$ & Area of one triangular finite element & m$^2$ \\
$N_a$ & Shape function associated with local node $a$ & dimensionless \\
$A_e$ & Area of element $e$ & m$^2$ \\
$\bm{C}$ & Vector of nodal defect concentrations & m$^{-3}$ \\
$\bm{M}$ & FEM mass matrix for the transient term & m$^2$ \\
$\bm{K}_D$ & FEM diffusion matrix & m$^2$/s \\
$\bm{K}_R$ & FEM reaction or annealing matrix & m$^2$/s \\
$\Delta t$ & Time-step size & s \\
$S$ & Volumetric defect source term & m$^{-3}$ s$^{-1}$ \\
\bottomrule
\end{longtable}

\section{Purpose of Module 3 and its role in the project}

Module 1 determines where radiation deposits energy. That deposited energy is the initial physical insult, but it is not yet the same thing as long-term damage. The material consequences of a radiation event evolve afterward. Defects can diffuse, recombine, cluster, dissociate, or anneal. Module 3 is the part of the project that turns the initial damage map into an evolving defect landscape.

This module is central because many downstream device consequences depend not only on how much damage was created, but also on where that damage migrates, how quickly it disappears, and whether it becomes more complex. Module 2 uses the defect population to modify space charge. Module 4 uses temperature to change mobility and reaction rates. Module 5 uses the evolving defect landscape to alter recombination and trapping. Thus Module 3 is the kinetic heart of the entire framework.

\section{Conservation-law starting point}

\subsection{Defect concentration as a continuum field}

Let $C(x,y,t)$ denote the concentration of a defect species, in units of number per volume. Even though defects are discrete lattice-scale objects, at the device scale it is useful to represent them by a coarse-grained concentration field. This is the same continuum logic used throughout transport theory: a sufficiently large but locally small averaging volume contains many defects, so a smooth field description becomes meaningful.

\subsection{Balance law for a defect species}

Consider a small material element. The concentration inside that element can change for three reasons:
\begin{enumerate}[label=\arabic*.]
    \item defects can flow into or out of the element,
    \item defects can be created inside the element,
    \item defects can be destroyed or transformed inside the element.
\end{enumerate}

The local conservation or balance law is therefore
\begin{equation}
\pdv{C}{t} + \nabla \cdot \mathbf{J}_C = R(C,T,\ldots),
\end{equation}
where $\mathbf{J}_C$ is the defect flux and $R$ is the local net reaction term. The right-hand side is positive when reactions create the species and negative when reactions consume it.

This equation is the natural defect analogue of the continuity equation. The key modeling task is therefore to specify the flux law and the reaction law.

\section{Diffusive transport from Fick's law}

\subsection{Why diffusion appears}

If the defect concentration is spatially nonuniform, there is a thermodynamic tendency for the distribution to smooth out. Regions of high concentration tend to feed regions of low concentration through random migration. At the continuum level, the simplest constitutive relation for that process is Fick's law:
\begin{equation}
\mathbf{J}_C = -D \nabla C.
\end{equation}
The minus sign is crucial. It says the flux points down the concentration gradient. Defects move from more concentrated regions toward less concentrated regions.

\subsection{From Fick's law to the diffusion equation}

Substituting Fick's law into the balance law gives
\begin{equation}
\pdv{C}{t} + \nabla \cdot (-D \nabla C) = R,
\end{equation}
so
\begin{equation}
\pdv{C}{t} = \nabla \cdot (D \nabla C) + R.
\end{equation}
If $D$ is spatially uniform, this becomes
\begin{equation}
\pdv{C}{t} = D \nabla^2 C + R.
\end{equation}
This is the diffusion-reaction equation. It is the core governing equation for Module 3.

\subsection{Physical meaning of the divergence form}

The divergence form,
\begin{equation}
\pdv{C}{t} = \nabla \cdot (D \nabla C) + R,
\end{equation}
is more fundamental than the constant-$D$ Laplacian form because it remains correct even when $D$ varies in space or depends on temperature. The diffusion operator is not merely ``curvature.'' It is the divergence of a flux. That distinction becomes essential once Module 4 temperature fields are coupled in and the diffusivity varies from place to place.

\section{Reaction and annealing terms}

\subsection{Why a reaction term is needed}

Radiation damage is not frozen in place. A vacancy may meet an interstitial and recombine. A mobile point defect may be absorbed at a sink. A defect complex may form from simpler constituents. Some defects may dissociate. These processes are not represented by diffusion alone. They require a local source-sink term.

Therefore the reaction term $R$ is not optional. It is the place where defect chemistry and defect kinetics enter the model.

\subsection{First-order annealing as the simplest reduced model}

The simplest physically meaningful reaction model is first-order annealing,
\begin{equation}
R = -k_{\mathrm{ann}} C,
\end{equation}
where $k_{\mathrm{ann}}$ is the annealing rate constant. The governing equation then becomes
\begin{equation}
\pdv{C}{t} = \nabla \cdot (D \nabla C) - k_{\mathrm{ann}} C.
\end{equation}
This says that defects spread by diffusion while being removed locally at a rate proportional to the current concentration.

\subsection{Why first-order annealing is useful}

First-order annealing is not the most general reaction law, but it is a highly useful reduced model because:
\begin{enumerate}[label=\arabic*.]
    \item it gives a clean analytic limit when diffusion is absent,
    \item it captures the idea of an intrinsic decay timescale,
    \item it is easy to verify numerically,
    \item it provides a baseline before introducing more complicated coupled kinetics.
\end{enumerate}

\subsection{General multi-species reaction structure}

For multiple defect species $C_i$, the balance law becomes
\begin{equation}
\pdv{C_i}{t} = \nabla \cdot (D_i \nabla C_i) + R_i(C_1,C_2,\ldots,T).
\end{equation}
This is the more general architecture the project can grow into. In the earliest implementation, however, a single effective defect species is sufficient to establish the transport framework.

\section{Reduction to the two-dimensional project form}

The project has adopted a 2D architecture for Modules 2--6. Thus the defect concentration becomes
\begin{equation}
C = C(x,y,t).
\end{equation}
For Cartesian coordinates,
\begin{equation}
\nabla \cdot (D \nabla C)
= \pdv{}{x}\left(D\pdv{C}{x}\right)
+ \pdv{}{y}\left(D\pdv{C}{y}\right).
\end{equation}
Therefore the 2D governing equation is
\begin{equation}
\pdv{C}{t}
=
\pdv{}{x}\left(D\pdv{C}{x}\right)
+
\pdv{}{y}\left(D\pdv{C}{y}\right)
+ R(C,T,\ldots).
\end{equation}
If the diffusivity is uniform and the reaction is first-order annealing,
\begin{equation}
\pdv{C}{t}
=
D\left(\pdv[2]{C}{x} + \pdv[2]{C}{y}\right) - k_{\mathrm{ann}} C.
\end{equation}
This is the baseline project equation for Module 3.

\section{Initial conditions from radiation damage maps}

\subsection{How Module 1 feeds Module 3}

Geant4 does not directly output an already-evolved defect concentration. It outputs deposited energy, event locations, and possibly a binned damage map. Module 3 must translate that radiation information into an initial defect condition or source term.

The simplest reduced mapping is
\begin{equation}
C(x,y,0) = C_0(x,y),
\end{equation}
where $C_0$ is obtained from a damage-generation rule based on the Geant4 deposition map. In a time-resolved variant one could instead use a source term $S(x,y,t)$, but for a first post-irradiation evolution model it is enough to treat radiation as an initial condition.

\subsection{Why this distinction matters}

This point is conceptually important: Geant4 determines where the initial insult occurs, while Module 3 determines what that insult becomes afterward. Conflating these two would erase the very defect-kinetics physics the project is trying to capture.

\section{Boundary conditions and their physical meaning}

\subsection{Zero-flux boundary condition}

The most natural default condition for a closed region with no defect exchange across the boundary is
\begin{equation}
\hat{\mathbf{n}} \cdot \mathbf{J}_C = 0.
\end{equation}
Using Fick's law, this becomes
\begin{equation}
\hat{\mathbf{n}} \cdot \nabla C = 0.
\end{equation}
This means there is no net defect flow across the boundary. Physically, the boundary acts as a reflecting boundary or symmetry plane.

\subsection{Dirichlet boundary condition}

Another possibility is to prescribe the concentration directly:
\begin{equation}
C = C_b \quad \text{on a boundary.}
\end{equation}
This is less common for a closed silicon device region but can represent contact with a reservoir or an imposed sink environment.

\subsection{Reactive or sink boundary condition}

A more physical future refinement is a Robin-type condition that represents surface absorption or sink behavior,
\begin{equation}
-\hat{\mathbf{n}}\cdot D\nabla C = h_s (C - C_{\mathrm{eq}}).
\end{equation}
This says the boundary can absorb or emit defects depending on deviation from an equilibrium value. Such a condition is useful when surfaces or interfaces act as strong sinks for mobile defects.

\section{Analytic limiting cases used for verification}

\subsection{Pure annealing with no diffusion}

If diffusion is switched off, $D=0$, and the reaction is first-order,
\begin{equation}
\pdv{C}{t} = -k_{\mathrm{ann}} C.
\end{equation}
The solution is
\begin{equation}
C(t) = C(0)e^{-k_{\mathrm{ann}} t}.
\end{equation}
If the initial state is spatially uniform, the entire field should remain spatially uniform while decaying exponentially. This is one of the cleanest verification tests because the exact solution is known.

\subsection{Pure diffusion with no reaction}

If $R=0$, the equation becomes
\begin{equation}
\pdv{C}{t} = D\nabla^2 C.
\end{equation}
Under zero-flux boundaries, the total amount of defect,
\begin{equation}
M(t) = \int C\,\dd V,
\end{equation}
should be conserved. The concentration profile should smooth and broaden, but the integral should remain constant.

\subsection{Uniform-state preservation}

If the initial condition is exactly uniform and the boundaries are consistent with that uniformity, then diffusion should generate no structure. A correct numerical scheme should preserve the uniform state up to roundoff.

\section{Characteristic timescales and regime interpretation}

\subsection{Diffusion timescale}

Let $L$ be a characteristic length scale. The diffusive relaxation time is approximately
\begin{equation}
\tau_D \sim \frac{L^2}{D}.
\end{equation}
This says diffusion becomes slower for larger structures and faster for larger diffusivity.

\subsection{Annealing timescale}

For first-order annealing, the intrinsic decay time is
\begin{equation}
\tau_{\mathrm{ann}} \sim \frac{1}{k_{\mathrm{ann}}}.
\end{equation}

\subsection{Competition between diffusion and annealing}

The ratio
\begin{equation}
\frac{\tau_D}{\tau_{\mathrm{ann}}} \sim \frac{L^2 k_{\mathrm{ann}}}{D}
\end{equation}
indicates which process dominates.
\begin{itemize}
    \item If $\tau_D \ll \tau_{\mathrm{ann}}$, diffusion redistributes defects before much annealing occurs.
    \item If $\tau_D \gg \tau_{\mathrm{ann}}$, defects disappear locally before they can migrate far.
\end{itemize}
This competition is central to the interpretation of results. It determines whether damage remains localized, spreads out, or fades quickly.

\section{Temperature coupling and Arrhenius physics}

A major reason Module 4 exists is that both diffusion and annealing are temperature sensitive. The compact Arrhenius formulas used in Module 3 are not arbitrary fitting expressions. They come from the idea that a thermally activated process must overcome an energy barrier before it can proceed.

\subsection{Why thermal activation produces an exponential}

Suppose a defect process requires access to a higher-energy configuration separated from the initial state by an energy barrier $E_a$. In thermal equilibrium, the probability of finding a system with energy $E$ is weighted by the Boltzmann factor
\begin{equation}
P(E) \propto \exp\left(-\frac{E}{k_B T}\right).
\end{equation}
Therefore, the fraction of defects with enough thermal energy to reach a transition state at energy barrier $E_a$ scales like
\begin{equation}
P(\text{barrier crossed}) \propto \exp\left(-\frac{E_a}{k_B T}\right).
\end{equation}
If the process also has an attempt frequency or attempt rate $\nu_0$, then the resulting rate takes the generic activated form
\begin{equation}
\Gamma(T) = \nu_0 \exp\left(-\frac{E_a}{k_B T}\right).
\end{equation}
This is the physical origin of Arrhenius behavior. The exponential is the thermal penalty for climbing an energy barrier, while the prefactor measures how often the system attempts to cross that barrier.

\subsection{Derivation of the diffusion Arrhenius law}

For point-defect migration in a crystal, diffusion can be viewed as a sequence of thermally activated jumps between neighboring lattice sites. Let $a$ be a characteristic jump distance and let $\Gamma_D$ be the jump frequency. For an unbiased random walk, the diffusivity scales as jump length squared times jump frequency. Up to an order-unity geometric factor,
\begin{equation}
D \sim a^2 \Gamma_D.
\end{equation}
More explicitly, for a simple random walk in $d$ spatial dimensions,
\begin{equation}
D = \frac{a^2 \Gamma_D}{2d}.
\end{equation}
If each jump requires the defect to overcome a migration barrier $E_D$, then the jump frequency is thermally activated:
\begin{equation}
\Gamma_D(T) = \nu_D \exp\left(-\frac{E_D}{k_B T}\right),
\end{equation}
where $\nu_D$ is an attempt frequency associated with lattice vibrations or local configurational attempts. Substituting this into the random-walk expression gives
\begin{equation}
D(T) = \frac{a^2 \nu_D}{2d} \exp\left(-\frac{E_D}{k_B T}\right).
\end{equation}
The prefactor is commonly grouped into a single constant,
\begin{equation}
D_0 \equiv \frac{a^2 \nu_D}{2d},
\end{equation}
so that
\begin{equation}
D(T) = D_0 \exp\left(-\frac{E_D}{k_B T}\right).
\end{equation}
This is the diffusion Arrhenius law used in Module 3.

The quantity $E_D$ is therefore not just a symbol inserted by hand. It is the effective \emph{migration activation energy}: the energy barrier that a defect must overcome to move from one metastable lattice configuration to another. In a more microscopic picture, $E_D$ may contain contributions from the local crystal environment, strain, charge state, and nearby impurities or defects. In the reduced continuum model, all of that microscopic complexity is compressed into the single effective barrier $E_D$ and prefactor $D_0$.

\subsection{Derivation of the annealing Arrhenius law}

Annealing is also a thermally activated process, but its physical meaning differs from diffusion. Annealing represents the disappearance or transformation of a defect population through processes such as recombination, structural rearrangement, dissociation of an unstable complex, or conversion into a lower-energy configuration.

Suppose a defect-removal channel requires crossing an activation barrier $E_{\mathrm{ann}}$. Then the microscopic rate per defect takes the same activated form,
\begin{equation}
\Gamma_{\mathrm{ann}}(T) = \nu_{\mathrm{ann}} \exp\left(-\frac{E_{\mathrm{ann}}}{k_B T}\right).
\end{equation}
If the annealing channel is first order, the number of defects removed per unit volume per unit time is proportional to the local concentration:
\begin{equation}
R_{\mathrm{ann}} = -\Gamma_{\mathrm{ann}}(T) \, C.
\end{equation}
We then identify the first-order annealing constant as
\begin{equation}
k_{\mathrm{ann}}(T) \equiv \Gamma_{\mathrm{ann}}(T) = k_0 \exp\left(-\frac{E_{\mathrm{ann}}}{k_B T}\right),
\end{equation}
with
\begin{equation}
k_0 \equiv \nu_{\mathrm{ann}}.
\end{equation}
Thus the first-order decay law
\begin{equation}
\pdv{C}{t} = -k_{\mathrm{ann}}(T) C
\end{equation}
comes directly from the assumption that each defect has an independent thermally activated probability per unit time of being removed.

The quantity $E_{\mathrm{ann}}$ is the effective \emph{annealing activation energy}. Physically, it is the barrier associated with the dominant defect-removal or defect-transformation pathway being represented in the reduced model. It need not be identical to $E_D$, because migration and annealing are not the same microscopic event. A defect may move relatively easily yet require a larger barrier to recombine, or vice versa.

\subsection{Why the prefactors and activation energies are called effective}

In real irradiated silicon, there is usually not just one defect species and not just one microscopic pathway. Different vacancies, interstitials, complexes, and charge states can each have different migration barriers and reaction barriers. If the continuum model uses a single species concentration $C$ and a single pair of Arrhenius parameters, then
\begin{equation}
D_0,\; E_D,\; k_0,\; E_{\mathrm{ann}}
\end{equation}
should be interpreted as \emph{effective parameters} summarizing the dominant kinetics over the temperature and timescale range of interest. They may come from atomistic calculations, experiments, or calibration to measured annealing curves.

\section{Microscopic origin of defect kinetic coefficients}

The simple Arrhenius forms used above can create the false impression that the defect diffusivity and annealing coefficient depend only on temperature. That is not physically correct. Temperature controls the probability of surmounting an activation barrier, but the barrier itself and the prefactor multiplying the Arrhenius exponential come from the material-specific energy landscape seen by the defect.

\subsection{Quantities that contribute to the effective diffusivity and annealing coefficient}

For the present project, the full physics picture behind $D$ and $k_{\mathrm{ann}}$ includes the following contributors:
\begin{enumerate}[label=\arabic*.]
    \item \textbf{Interatomic bonding.} Defect motion requires stretching, weakening, or breaking local bonds as the defect passes from one metastable site to another. Bond strengths therefore contribute directly to the migration barrier and to the annealing barrier.
    \item \textbf{Local lattice geometry.} The crystal structure determines which sites are available, the jump distance between sites, the coordination around the defect, and the shape of the saddle-point pathway. This influences both the activation barrier and the diffusion prefactor.
    \item \textbf{Electronic structure.} Defect energies depend on orbital hybridization, local electron density, screening, and the location of defect levels relative to the bands. This changes the energies of stable and transition configurations and therefore changes both $E_D$ and $E_{\mathrm{ann}}$.
    \item \textbf{Defect charge state.} A charged defect can have a different migration barrier and a different reaction pathway from the neutral version of the same structural defect. Charge state is therefore one of the most important state variables once semiconductor electrostatics are considered.
    \item \textbf{Strain field.} Strain shifts the energies of metastable configurations and transition states. In a damaged region, near an interface, or in a mechanically stressed device, strain can bias defect motion and modify annealing kinetics.
    \item \textbf{Nearby impurities and dopants.} Dopants and impurities alter the local environment through binding interactions, electrostatic effects, and Fermi-level shifts. They can trap defects, stabilize some charge states, or lower or raise migration barriers.
    \item \textbf{Nearby interfaces or sinks.} Surfaces, oxide interfaces, and material boundaries can act as sinks or preferred defect environments. They therefore alter both the transport problem and the effective local reaction pathways.
    \item \textbf{Local electric field.} If the relevant defect is charged, an electric field can bias defect motion and can also lower or raise an activation barrier through electrostatic work along the migration path.
    \item \textbf{Local temperature.} Temperature controls the thermal activation probability and can also change phonon spectra, charge-state populations, and the relative importance of competing reaction pathways.
\end{enumerate}

The important point is that temperature is only one part of the story. The reduced continuum coefficients are temperature activated, but their detailed values arise from a much richer material-specific defect physics.

\subsection{Which of these belong directly in Module 3?}

Not all of the preceding physics should be promoted to independent continuum fields in the first implementation. A useful separation is:
\begin{itemize}
    \item \textbf{Microscopic origin level:} interatomic bonding, local lattice geometry, and electronic structure. These are the upstream material-physics ingredients that determine the barrier heights, attempt frequencies, and charge-state dependence.
    \item \textbf{Simulation-level modifier fields:} temperature, defect charge state, dopant environment, electric field, strain, and interface proximity. These are the variables that can realistically be inserted into a reduced continuum coefficient model.
\end{itemize}

That separation keeps the project honest. Module~3 does not solve orbital-scale or bond-scale physics directly. Instead, it uses reduced constitutive laws whose coefficients are understood to have been shaped by that microscopic physics.

\subsection{Reduced constitutive hierarchy adopted in this project}

The long-term constitutive goal is not merely
\begin{equation}
D = D(T), \qquad k_{\mathrm{ann}} = k_{\mathrm{ann}}(T),
\end{equation}
but rather
\begin{equation}
D = D\left(T, q_d, N_{\mathrm{dop}}, |\mathbf{E}|, \varepsilon, \xi_{\mathrm{int}}, \ldots\right),
\end{equation}
and
\begin{equation}
k_{\mathrm{ann}} = k_{\mathrm{ann}}\left(T, q_d, N_{\mathrm{dop}}, |\mathbf{E}|, \varepsilon, \xi_{\mathrm{int}}, \ldots\right).
\end{equation}
Here $q_d$ is the effective defect charge state, $N_{\mathrm{dop}}$ is a dopant-environment measure, $|\mathbf{E}|$ is the electric-field magnitude, $\varepsilon$ is a reduced strain measure, and $\xi_{\mathrm{int}}$ is an interface-proximity variable.

\subsection{What is implemented now}

The first implemented reduced model keeps the hierarchy manageable by activating only the dependencies that are expected to matter most at the current project stage:
\begin{equation}
D = D\left(T, q_d, N_{\mathrm{dop}}, |\mathbf{E}|\right),
\end{equation}
and
\begin{equation}
k_{\mathrm{ann}} = k_{\mathrm{ann}}\left(T, q_d, N_{\mathrm{dop}}, |\mathbf{E}|\right).
\end{equation}
Strain and interface proximity are explicitly documented as next-pass refinements but are not turned on in the first numerical implementation. Interatomic bonding, local lattice geometry, and electronic structure are treated as the microscopic origin of the coefficient parameters rather than as direct continuum fields solved in the MATLAB code.

\subsection{One reduced form consistent with the physics}

A practical reduced constitutive form that preserves the Arrhenius structure is
\begin{equation}
D = D_0^{\mathrm{eff}}(q_d, N_{\mathrm{dop}})
\exp\left[-\frac{E_D^{\mathrm{eff}}(q_d, N_{\mathrm{dop}}) - \Delta E_{\mathrm{field}}(|\mathbf{E}|)}{k_B T}\right],
\end{equation}
and
\begin{equation}
k_{\mathrm{ann}} = k_0^{\mathrm{eff}}(q_d, N_{\mathrm{dop}})
\exp\left[-\frac{E_{\mathrm{ann}}^{\mathrm{eff}}(q_d, N_{\mathrm{dop}}) - \Delta E_{\mathrm{field}}(|\mathbf{E}|)}{k_B T}\right].
\end{equation}
This is still a reduced continuum model, but it now makes clear which physics is being compressed into effective barriers and prefactors and which state variables are actively modifying those coefficients during the simulation.

\subsection{How Module 2 enters and how it does not}

It is important not to confuse the microscopic electric potential associated with bonding and orbital structure with the device-scale electrostatic potential solved in Module~2. Module~2 computes the continuum electric potential and electric field of the device cross section. That field can influence charged-defect transport or lower an activation barrier, so Module~2 can feed into Module~3 through $|\mathbf{E}|$. However, the bond-scale and orbital-scale origin of $E_D$ and $E_{\mathrm{ann}}$ is not something Module~2 itself computes. That physics belongs to the microscopic origin of the reduced coefficients.

\subsection{Module 3 equation with thermal coupling}

Once Module 4 supplies $T(x,y,t)$, the Module 3 equation becomes
\begin{equation}
\pdv{C}{t}
=
\nabla \cdot \left(D(T)\nabla C\right) - k_{\mathrm{ann}}(T)C.
\end{equation}
If $D$ varies spatially through its temperature dependence, then the diffusion term can be expanded as
\begin{equation}
\nabla \cdot \left(D(T)\nabla C\right) = D(T)\nabla^2 C + \nabla D(T) \cdot \nabla C.
\end{equation}
The first term is the familiar Laplacian smoothing term. The second term appears because temperature gradients create spatial gradients in diffusivity, which in turn bias how quickly different regions relax.

In 2D Cartesian form, the coupled equation is
\begin{equation}
\pdv{C}{t}
=
\pdv{}{x}\left(D(T)\pdv{C}{x}\right)
+
\pdv{}{y}\left(D(T)\pdv{C}{y}\right)
-
k_{\mathrm{ann}}(T)C.
\end{equation}
Now the thermal field changes not only how quickly defects move, but also how quickly they disappear. Higher temperature generally increases both diffusivity and annealing rate because the Boltzmann penalty for barrier crossing becomes less severe.

This coupling is one of the strongest reasons the project must be multiphysics. Defect kinetics and temperature cannot be treated as unrelated after the architecture has matured.


\section{Finite-element formulation of Module 3}

The preceding sections define the continuum physics. For implementation and interview defense, it is equally important to explain how the time-dependent diffusion-reaction equation is converted into algebraic equations. The present project now keeps two numerical paths conceptually distinct. The older Module 3 MATLAB path advances the equation on a structured Cartesian grid with a conservative finite-volume/finite-difference stencil. The new FEM path described here treats the same governing equation with linear triangular finite elements. Both solve the same physics, but they approximate the spatial operators differently.

\subsection{Strong form used for the FEM derivation}

Let $\Omega$ be the two-dimensional silicon cross section and let $\partial\Omega$ be its boundary. The scalar unknown is the defect concentration $C(x,y,t)$. The strong form is
\begin{equation}
\pdv{C}{t}
=
\nabla\cdot\left(D\nabla C\right)-k_{\mathrm{ann}}C+S
\qquad \text{in } \Omega,
\end{equation}
where $D$ is the defect diffusivity, $k_{\mathrm{ann}}$ is the first-order annealing coefficient, and $S$ is an optional defect source term. The default post-irradiation cases use $S=0$ after the initial defect field has been created.

The zero-flux boundary condition is
\begin{equation}
\hat{\mathbf{n}}\cdot\mathbf{J}_C = 0,
\qquad
\mathbf{J}_C=-D\nabla C,
\end{equation}
so
\begin{equation}
\hat{\mathbf{n}}\cdot D\nabla C = 0
\qquad \text{on the zero-flux part of } \partial\Omega.
\end{equation}
This is the natural default condition for a closed computational region. It says defects do not leave the modeled silicon region through the boundary.

\subsection{Weak form and why it is needed}

To derive the FEM equations, multiply the strong form by a test function $w(x,y)$ and integrate over the domain:
\begin{equation}
\int_{\Omega} w \pdv{C}{t}\,d\Omega
=
\int_{\Omega} w\nabla\cdot(D\nabla C)\,d\Omega
-
\int_{\Omega} wk_{\mathrm{ann}}C\,d\Omega
+
\int_{\Omega} wS\,d\Omega.
\end{equation}
Apply integration by parts to the diffusion term:
\begin{equation}
\int_{\Omega} w\nabla\cdot(D\nabla C)\,d\Omega
=
\int_{\partial\Omega} wD\nabla C\cdot\hat{\mathbf{n}}\,d\Gamma
-
\int_{\Omega} D\nabla w\cdot\nabla C\,d\Omega.
\end{equation}
For a homogeneous zero-flux boundary, the boundary integral vanishes. The weak form becomes
\begin{equation}
\boxed{
\int_{\Omega} w \pdv{C}{t}\,d\Omega
+
\int_{\Omega} D\nabla w\cdot\nabla C\,d\Omega
+
\int_{\Omega} wk_{\mathrm{ann}}C\,d\Omega
=
\int_{\Omega} wS\,d\Omega .
}
\end{equation}
This equation is the FEM starting point. It contains only first spatial derivatives of $C$, which is why ordinary $C^0$ continuous linear triangular elements are sufficient. In physical terms, the weak form balances the rate of defect inventory change, diffusive smoothing, local annealing loss, and any volumetric source.

\subsection{Triangular mesh and nodal unknowns}

The FEM implementation represents $\Omega$ using a triangular mesh. Each triangular element has three local nodes. The global unknown vector is
\begin{equation}
\bm{C}(t)=
\begin{bmatrix}
C_1(t) & C_2(t) & \cdots & C_N(t)
\end{bmatrix}^{T},
\end{equation}
where $C_A(t)$ is the defect concentration at global node $A$ and $N$ is the total number of mesh nodes. This is analogous to the electrostatic potential vector in Module 2, except the unknown is now concentration rather than potential and the equation is time dependent rather than steady.

A rectangular silicon cross section can still be meshed with triangles by first making a structured grid of points and then splitting each rectangular cell into two counter-clockwise triangular elements. This gives a simple mesh whose connectivity is easy to inspect and whose boundary node sets are known exactly.

\subsection{Linear triangular shape functions}

Inside element $e$, the concentration is approximated by
\begin{equation}
C^{(e)}(x,y,t)
\approx
\sum_{a=1}^{3}N_a(x,y)C_a^{(e)}(t),
\end{equation}
where $N_a$ are the three linear triangular shape functions. Each shape function is one at its own local node and zero at the other two local nodes:
\begin{equation}
N_a(\xvec_b)=\delta_{ab}.
\end{equation}
Thus the nodal values define the concentration everywhere inside the element by interpolation.

For a linear triangle,
\begin{equation}
N_a(x,y)=\alpha_a+\beta_a x+\gamma_a y,
\end{equation}
so the gradient $\nabla N_a$ is constant within the element. This makes the element diffusion matrix especially simple, because $\nabla C$ is also constant inside each linear triangle.

\subsection{Element mass matrix}

The first term in the weak form produces a mass matrix. It is called a mass matrix by analogy with transient structural dynamics, but here it does not represent mechanical mass. It represents the storage or inventory of the scalar concentration field. The element mass matrix is
\begin{equation}
M_{ab}^{(e)}=\int_{\Omega_e}N_aN_b\,d\Omega.
\end{equation}
For a linear triangular element with area $A_e$,
\begin{equation}
\bm{M}^{(e)}=
\frac{A_e}{12}
\begin{bmatrix}
2 & 1 & 1 \\
1 & 2 & 1 \\
1 & 1 & 2
\end{bmatrix}.
\end{equation}
This matrix converts nodal time derivatives into element-level inventory changes.

\subsection{Element diffusion matrix}

The diffusion term gives
\begin{equation}
K_{D,ab}^{(e)}=\int_{\Omega_e}D\nabla N_a\cdot\nabla N_b\,d\Omega.
\end{equation}
If $D$ is constant over the element,
\begin{equation}
K_{D,ab}^{(e)}=D A_e\nabla N_a\cdot\nabla N_b.
\end{equation}
This matrix penalizes sharp concentration gradients between neighboring nodes. It is the defect-transport analogue of the electrostatic stiffness matrix in Module 2. The word ``stiffness'' here is numerical, not mechanical: it quantifies how strongly diffusion resists spatial concentration variations.

When $D$ is supplied as a spatial field, the first implementation uses an element-average diffusivity,
\begin{equation}
D_e \approx \frac{D_1+D_2+D_3}{3},
\end{equation}
where $D_1,D_2,D_3$ are the nodal values on that triangle. Higher-order quadrature can be added later if the diffusivity varies strongly within individual elements.

\subsection{Element reaction or annealing matrix}

The first-order annealing term gives
\begin{equation}
K_{R,ab}^{(e)}=\int_{\Omega_e}k_{\mathrm{ann}}N_aN_b\,d\Omega.
\end{equation}
If $k_{\mathrm{ann}}$ is constant within the element,
\begin{equation}
\bm{K}_{R}^{(e)}=k_{\mathrm{ann},e}\bm{M}^{(e)}.
\end{equation}
This is the finite-element version of the local loss term $-k_{\mathrm{ann}}C$. It removes concentration locally without requiring spatial gradients. If the annealing coefficient varies in space, the first implementation again uses an element-average value.

\subsection{Element source vector}

If a volumetric source $S(x,y,t)$ is active, the element source vector is
\begin{equation}
f_a^{(e)}=\int_{\Omega_e}N_aS\,d\Omega.
\end{equation}
For the default post-irradiation evolution problem, radiation damage is represented mainly through the initial condition $C(x,y,0)=C_0(x,y)$ rather than a continuing source. In that case $\bm{f}^{(e)}=\bm{0}$ after $t=0$. If a time-dependent radiation source is later introduced, this vector is the natural place where it enters the FEM formulation.

\subsection{Global assembly}

Each triangular element contributes local matrices $\bm{M}^{(e)}$, $\bm{K}_{D}^{(e)}$, $\bm{K}_{R}^{(e)}$, and a local source vector $\bm{f}^{(e)}$. Assembly adds these local contributions into global matrices using the element-to-global node connectivity:
\begin{align}
M_{AB} &\leftarrow M_{AB}+M_{ab}^{(e)},\\
K_{D,AB} &\leftarrow K_{D,AB}+K_{D,ab}^{(e)},\\
K_{R,AB} &\leftarrow K_{R,AB}+K_{R,ab}^{(e)},\\
f_A &\leftarrow f_A+f_a^{(e)}.
\end{align}
Here local node $a$ corresponds to global node $A$, and local node $b$ corresponds to global node $B$.

The semi-discrete FEM equation is therefore
\begin{equation}
\boxed{
\bm{M}\dv{\bm{C}}{t}+\left(\bm{K}_D+\bm{K}_R\right)\bm{C}=\bm{f}.
}
\end{equation}
This is the central FEM equation for Module 3. It should be read as the matrix form of conservation plus diffusion plus annealing.

\subsection{Implicit time stepping}

The first FEM implementation uses backward Euler time stepping because it is more stable for diffusion-reaction problems than a simple explicit update. With time step $\Delta t$, the semi-discrete equation becomes
\begin{equation}
\left[\bm{M}+\Delta t\left(\bm{K}_D+\bm{K}_R\right)\right]\bm{C}^{n+1}
=
\bm{M}\bm{C}^{n}+\Delta t\bm{f}^{n+1}.
\end{equation}
For source-free evolution,
\begin{equation}
\left[\bm{M}+\Delta t\left(\bm{K}_D+\bm{K}_R\right)\right]\bm{C}^{n+1}
=
\bm{M}\bm{C}^{n}.
\end{equation}
The main advantage is that the method avoids the severe diffusion stability limit that appears in explicit schemes. Time-step accuracy still matters, but the method is not constrained in the same way by $\Delta t \lesssim h^2/D$ for stability.

\subsection{Boundary-condition handling}

Homogeneous zero-flux boundaries are natural in the weak form. Because the boundary integral vanishes, no special matrix modification is required. This is different from a Dirichlet boundary, where the concentration is prescribed directly and the corresponding nodal values must be constrained.

For most Module 3 cases, zero-flux boundaries are the correct default because they test whether defect inventory is conserved under pure diffusion. Dirichlet or Robin sink boundaries can be added later to model surfaces, interfaces, or contacts that absorb mobile defects.

\subsection{Postprocessing and conserved quantities}

After solving for $\bm{C}^{n+1}$, the defect concentration can be plotted directly at nodes or interpolated over the triangular mesh. The defect flux inside an element is
\begin{equation}
\mathbf{J}_C^{(e)}=-D_e\nabla C^{(e)}
=-D_e\sum_{a=1}^{3}C_a^{(e)}\nabla N_a.
\end{equation}
For linear triangles this flux is piecewise constant over each element.

The total defect inventory is computed in FEM form as
\begin{equation}
M_C(t)=\int_{\Omega}C(x,y,t)\,d\Omega
\approx
\mathbf{1}^{T}\bm{M}\bm{C}(t),
\end{equation}
where $\mathbf{1}$ is a vector of ones. This quantity should remain constant for pure diffusion with zero-flux boundaries and should decay exponentially for spatially uniform first-order annealing.

\subsection{MATLAB implementation map}

The FEM implementation is organized as a parallel path to the existing structured-grid Module 3 solver:
\begin{longtable}{p{0.40\textwidth} p{0.50\textwidth}}
\toprule
\textbf{File} & \textbf{Role} \\
\midrule
\endhead
\path{main_module3_fem_defect_evolution.m} & Top-level driver for the FEM Module 3 cases. \\
\path{default_module3_fem_params.m} & Defines the FEM domain, mesh resolution, diffusivity, annealing rate, time step, and initial condition. \\
\path{make_rectangular_tri_mesh_2d.m} & Reuses the triangular mesh generator introduced for Module 2. \\
\path{initialize_defect_field_fem_2d.m} & Builds nodal initial defect concentration fields. \\
\path{compute_element_mass_triangle.m} & Computes the consistent mass matrix for a linear triangular element. \\
\path{compute_element_diffusion_triangle.m} & Computes the triangular diffusion matrix from $D\nabla N_a\cdot\nabla N_b$. \\
\path{assemble_defect_fem_2d.m} & Assembles global mass, diffusion, reaction, and source arrays. \\
\path{solve_defect_diffusion_reaction_fem_2d.m} & Advances the implicit FEM system through time and records diagnostics. \\
\path{write_module3_fem_summary.m} & Writes a text summary of the FEM run and verification metrics. \\
\bottomrule
\end{longtable}

\subsection{What an interviewer may ask about the FEM side}

A concise technical answer is: Module 3 solves a time-dependent scalar diffusion-reaction equation. The FEM weak form is obtained by multiplying by a test function, integrating over the domain, and integrating the diffusion term by parts. Linear triangular elements approximate the concentration field. The transient term produces a mass matrix $\bm{M}$, diffusion produces $\bm{K}_D$, first-order annealing produces $\bm{K}_R$, and backward Euler gives $[\bm{M}+\Delta t(\bm{K}_D+\bm{K}_R)]\bm{C}^{n+1}=\bm{M}\bm{C}^{n}+\Delta t\bm{f}^{n+1}$. Zero-flux boundaries are natural boundary conditions and require no explicit constraint in the homogeneous case.


\section{How Module 3 couples to other modules}

\subsection{Input from Module 1}

Module 1 supplies the initial damage map or source distribution.

\subsection{Output to Module 2}

Module 3 provides the evolving concentrations $C_i(x,y,t)$ that are converted into defect charge density in Module 2.

\subsection{Input from Module 4}

Module 4 provides the temperature field that changes diffusivities and reaction rates.

\subsection{Output to Module 5}

Module 5 uses the defect distribution to model trap-assisted recombination, trapping, and lifetime degradation.

Thus Module 3 is not isolated. It is the evolving material-state layer that communicates with all of the main downstream device-physics modules.

\section{What this module captures and what it does not}

Module 3 captures continuum-scale migration and reaction of defects. It does not directly resolve atomistic trajectories or explicitly compute formation energies from first principles. Those deeper physics can be folded into the continuum parameters $D$, $R$, and their temperature dependence, but the module itself is a device-scale defect-kinetics model rather than a molecular-dynamics or density-functional-theory solver.

This distinction is useful because it keeps the project scope realistic while still preserving physically meaningful links to microscopic input data.

\section{Summary of the final project equation}

The baseline 2D Module 3 equation is
\begin{equation}
\pdv{C}{t}
=
\pdv{}{x}\left(D\pdv{C}{x}\right)
+
\pdv{}{y}\left(D\pdv{C}{y}\right)
+ R(C,T,\ldots).
\end{equation}
For the first reduced implementation,
\begin{equation}
\pdv{C}{t}
=
D\left(\pdv[2]{C}{x} + \pdv[2]{C}{y}\right) - k_{\mathrm{ann}}C.
\end{equation}
With temperature coupling,
\begin{equation}
\pdv{C}{t}
=
\nabla\cdot\left(D(T)\nabla C\right) - k_{\mathrm{ann}}(T)C.
\end{equation}
This is the equation set that turns a radiation-created damage map into an evolving defect landscape suitable for coupling to electrostatics, thermal physics, and carrier transport.

\section{References}
\begin{sloppypar}
\begin{enumerate}[label={[\arabic*]}]
\item A. Fick, ``Ueber Diffusion, \emph{Annalen der Physik} \textbf{170}, 59--86 (1855).
\item J. Crank, \emph{The Mathematics of Diffusion}, 2nd ed., Oxford University Press (1975).
\item P. M. Fahey, P. B. Griffin, and J. D. Plummer, ``Point defects and dopant diffusion in silicon, \emph{Reviews of Modern Physics} \textbf{61}, 289--384 (1989).
\item G. D. Watkins, ``Intrinsic defects in silicon, in \emph{Deep Centers in Semiconductors}, 2nd ed., Gordon and Breach (1999).
\item S. M. Sze and K. K. Ng, \emph{Physics of Semiconductor Devices}, 3rd ed., Wiley (2007).
\item O. C. Zienkiewicz, R. L. Taylor, and J. Z. Zhu, \emph{The Finite Element Method: Its Basis and Fundamentals}, 7th ed., Elsevier (2013).
\item D. L. Logan, \emph{A First Course in the Finite Element Method}, 6th ed., Cengage (2017).
\end{enumerate}
\end{sloppypar}

\end{document}
